\documentclass[11pt]{article}
\input{packages.tex}
\usepackage{soul}
\DeclareMathOperator{\sign}{sign}
\DeclareMathOperator{\Imm}{Im}
\DeclareMathOperator{\rg}{rg}

% Настройка колонтитулов
\pagestyle{fancy}
\fancyhf{} % Очистить все поля
\fancyhead[L]{Ура! Победа нас настигла!}
\fancyhead[R]{\href{https://t.me/lost_luna}{Автор}}
\fancyfoot[C]{\thepage}               % Справа внизу

\setlength{\headheight}{13.6pt}

\hypersetup{
    colorlinks=true,
    linkcolor=blue,
    citecolor=blue,
    filecolor=blue,
    urlcolor=blue,
}

\begin{document}

\begin{titlepage}
    \centering
    \vspace*{\fill}  % заполняет пространство сверху

    {\Huge Лекция 7. Характеристическое уравнение. Свойства собственных подпространств.
    Комплексные характеристические числа. Диагонализируемое преобразование. 
    \par}

    \vspace*{\fill}  % заполняет пространство снизу
\end{titlepage}

\setcounter{page}{1}

\section{Характеристическое уравнение}

\begin{theorem}
    Пусть линейное преобразование $\varphi$ имеет матрицы преобразования $A$ и $A'$
    в разных базисах, тогда характеристические многочлены этих матриц совпадают.
\end{theorem}

\begin{proof}
    Из замечательной пятой лекции нам известно, что \[A' = S^{-1} \cdot A \cdot S\]
    Тогда рассмотрим \[\det(A' - \lambda E) = \det(S^{-1} \cdot A \cdot S - \lambda E) =
    \det(S^{-1} \cdot (A - \lambda E) \cdot S) = \det S^{-1} \cdot \det(A - \lambda E) \cdot \det S\]
    \[= [\det S \cdot \det S^{-1} = \det (S \cdot S^{-1}) = \det E = 1] = \det (A - \lambda E)\]
\end{proof}

\begin{consequence}
    Характеристический многочлен матрицы $A$ можно назвать характеристическим 
    многочленом \textbf{преобразования} $\varphi$.
\end{consequence}

\begin{definition}
    Следом матрицы $A$ называется коэффициент \[a_{11} + a_{22} + ... + a_{nn}
    \text{ при } (-\lambda)^{n - 1}\]
\end{definition}

\begin{consequence}
    Детерминант матрицы $A$ и коэффициенты характеристического многочлена (в частности,
    след матрицы $A$) являются инвариантами (то есть сохраняются в любом базисе).
\end{consequence}

\section{Свойства собственных подпространств}

\begin{theorem}
    Сумма собственных подпространств является прямой. \\ Эквивалентная формулировка:
    собственные векторы, отвечающие различным собственным значениям, линейно независимы.
\end{theorem}

\begin{proof}
    Пусть $h_1, \dots, h_S$ --- собственные вектора, отвечающие различным собственным 
    значениям $\lambda_1, \dots, \lambda_s$. Пусть $A$ --- матрица преобразования,
    относительно которого рассматриваются собственные подпространства. Рассмотрим 
    преобразование вида \[B_i = A - \lambda_i E, \quad i \in \{1, \dots, S\}\]
    Тогда для любых $i$ и $j$ имеем: \[B_i(h_j) = (A - \lambda_i E) \cdot h_j =
    A \cdot h_j - \lambda_i E \cdot h_j = \lambda_j \cdot h_j - \lambda_i \cdot h_j = 
    (\lambda_j - \lambda_i) \cdot h_j\] Следовательно, \[B_i(h_j) = 0 \Longleftrightarrow i = j\]
    Предположим, что $\{h_1, \dots, h_S\}$ --- линейно зависимый набор. Без ограничения 
    общности считаем, что \[h_1 = \alpha_2 \cdot h_2 + \dots + \alpha_S \cdot h_S \label{eq:star1} \tag{*}\]
    Последовательно подействуем на (\ref{eq:star1}) преобразованиями $B_2, \dots, B_S$
    и получим: \[(\lambda_1 - \lambda_2) \cdot (\lambda_1 - \lambda_3) \cdot ... \cdot 
    (\lambda_1 - \lambda_s) = 0 + ... + 0\]
    Так как в каждом слагаемом в правой части всегда найдется скобка вида $(\lambda_k - \lambda_k)$.
    Получим, что ненулевой вектор равен нулевому $\Longrightarrow$ противоречие.
\end{proof}

\begin{theorem}
    Пусть $\lambda_1$ --- корень характеристического многочлена кратности $s$ (то есть
    многочлен делится на $(\lambda - \lambda_1)^s$ без остатка). Тогда $\lambda_1$
    соответствует собственное подпространство $\mathcal{L}_1$, причем 
    \[1 \leq \dim \mathcal{L}_1 \leq s\]
\end{theorem}

\begin{proof}
    Предположим, что $\dim \mathcal{L}_1 = k$. Выберем базис в $\mathcal{L}$ так,
    что первые $k$ базисных векторов $e_1, \dots, e_k$ --- базис в $\mathcal{L}_1$,
    тогда матрица преобразования имеет следующий вид: 
    \[ A = \left(
    \begin{array}{c|c}
        \begin{array}{ccc}
        \lambda_1 & \dots & 0 \\ \vdots & \ddots & \vdots \\ 0 & \dots & \lambda_1
        \end{array}
        & \widetilde{A} \\ 0 
    \end{array}\right)\]
    Так как $A(e_i) = \lambda_1 \cdot e_1 \ \forall i \in \{1, \dots, k\}$. Раскладывая
    детерминант матрицы $A - \lambda E$ последовательно по каждому из первых $k$
    столбцов, получим \[\det (A - \lambda E) = (\lambda_1 - \lambda)^k \cdot Q(\lambda) = 
    [ \text{ но в силу тайного знания кратности } \lambda_1] = (\lambda_1 - \lambda)^s \cdot R(\lambda)\]
    где в многочлене $R(\lambda)$ нет множителей вида $(\lambda_1 - \lambda)$. 
    Следовательно, \[k \leq s\] а неравенство $1 \leq k$ тривиально выполняется.
\end{proof}

\begin{definition}
    Число $s$ будем называть \textbf{алгебраической кратностью}, а число $k$ ---
    \textbf{геометрической кратностью}.
\end{definition}

\begin{example}
    Рассмотрим преобразование, в котором собственному значению кратности $s$
    принадлежит собственное подпространство размерности меньшей, чем $s$:
    \[A = \begin{pmatrix}
        1 & 1 \\ 0 & 1
    \end{pmatrix}\]
    \[\det (A - \lambda E) = \det \begin{pmatrix}
        1 - \lambda & 1 \\ 0 & 1 - \lambda
    \end{pmatrix} = (1 - \lambda)^2 = 0\]
    Собственное значение $\lambda = 1$ имеет кратность $2$, но \[(A - \lambda E) \cdot h 
    = 0 \Longleftrightarrow \begin{pmatrix}
        0 & 1 \\ 0 & 0
    \end{pmatrix} \cdot h \Longleftrightarrow h = \begin{pmatrix}
        1 \\ 0
    \end{pmatrix} \text{ --- базис в собственнном подпространстве}\]
    Получили, что собственному значению кратности $2$ соответствует одномерное
    собственное подпространство.
\end{example}

\section{Комплексные характеристические числа}

\begin{remark}
    Пусть у линейного преобразования $A$ характеристический многочлен имеет
    комплексный корень $\lambda$, тогда комплексно сопряженное число $\overline{\lambda}$
    тоже будет корнем многочлена (поскольку у многочлена вещественные коэффициенты).
\end{remark}

\begin{theorem}
    Паре комплексно сопряженных корней характеристического многочлена преобразования
    $A$ соответствует ненулевое инвариантное подпространство $\mathcal{L}_1$ такое, что 
    \begin{enumerate}
        \item $\mathcal{L}_1$ не содержит собственных векторов 
        \item через любой вектор $\mathcal{L}_1$ проходит двумерное инвариантное подпространство
    \end{enumerate}
\end{theorem}

\begin{proof}
    Зафиксируем $\lambda$ и $\overline{\lambda}$, им соответствует трехчлен \[t^2 + p \cdot t + q\]
    корнями которого они являются. Рассмотрим преобразование \[B^2 = A^2 + p A + q E\]
    Пусть подпространство $\mathcal{L}_1 = Ker B$. Из теоремы в предыдущей лекции,
    $\mathcal{L}_1$ инвариантно относительно $A$, так как $A$ и $B$ перестановочны.
    Заметим, что \[ B = (A - \lambda E) \cdot (A - \overline{\lambda} E) \Longrightarrow
    \det B = \det (A - \lambda E) \cdot \det (A - \overline{\lambda} E) = 0\]
    Так как $\lambda$ --- корень характеристического многочлена, то
    $\det (A - \lambda E) = 0$. Из этого всего великолепия следует, что $B$ имеет 
    ненулевое ядро, то есть $\mathcal{L}_1$ --- ненулевое подпространство. \\
    Предположим, что \[\exists h \in \mathcal{L}_1: A \cdot h = \lambda_1 \cdot h\]
    то есть в пространстве $\mathcal{L}_1$ (или же в $Ker B$) имеется собственный вектор:
    \[h \in Ker B \Longrightarrow B \cdot h = (A^2 + p A + q E) \cdot h = 0 
    \Longrightarrow A^2 + p A + q E = 0 \ (\text{так как } h \ne 0)\]  
    Но уравнение $A^2 + p A + q E = $ имеет лишь комплексные корни, а $\lambda_1 \in \mathbb{R}$.
    Следовательно, получаем противоречие. \\
    Пусть теперь $x \ne 0$ и $x \in \mathcal{L}_1$. Рассмотрим $A(x)$. Покажем, что 
    $\mathcal{L}_2 = \ <x, \ A(x)>$ --- инвариантное подпространство. Рассмотрим 
    \[a = \alpha x + \beta A x\] \[A(a) = \alpha A x + \beta A^2 x \label{eq:star2} \tag{*}\]
    При этом $x \in Ker B$, следовательно, \[(A^2 + p A + q E) x = 0 \ \Longrightarrow
    \ A^2 x = -p A x - q E x\]
    Подставим это в (\ref{eq:star2}): \[A(a) = (\alpha - p) A x - q x \ \Longrightarrow
    \ A(a) \in \ <x, \ A(x)>\]
    Осталось показать, что $\mathcal{L}_2$ двумерно. Ясно, что $\dim \mathcal{L}_2 \leq 2$.
    При этом, если бы $x$ и $A(x)$ лежали в одномерном пространстве, то $x$ был бы собственным 
    вектором, а это противоречит доказанному пункту 1.
\end{proof}

\begin{consequence}
    У любого преобразования конечномерного ненулевого пространства всегда найдется 
    либо одномерное собственное подпространство, либо двумерное инвариантное.
\end{consequence}

\section{Диагонализируемое преобразование}

\begin{definition}
    Преобразование называется диагонализируемым, если существует базис, в котором 
    матрица преобразования имеет диагональный вид. \\ 
    Эквивалентная формулировка: существует базис из собственных векторов
\end{definition}

\begin{theorem}[Необходимое условие диагонализируемости]
    Если преобразование $A$ линейного пространства $\mathcal{L}$ диагонализируемо, 
    то все корни характеристического уравнения вещественны. 
\end{theorem}

\begin{proof}
    Если у характеристического многочлена есть комплексные корни, то он имеет вид:
    \[(\dots) \cdot (\dots) \cdot ... \cdot \underset{\lambda \in \mathbb{C}}{(\dots)}\]
    Если преобразование диагонализируемо, то существует базис, в котором матрица преобразования 
    имеет диагональный вид, а значит, характеристический многочлен имеет вид: 
    \[(\lambda - \lambda_1) \cdot ... \cdot (\lambda - \lambda_n), \quad \lambda_i \in  
    \mathbb{R} \ \forall i \in \{1, \dots, n\}\]
    Но характеристический многочлен не зависит от базиса, следовательно, получаем противоречие.
\end{proof}

\begin{theorem}[Достаточное условие диагонализируемости]
    Преобразование $A$ линейного пространства $\mathcal{L}$ диагонализируемо, если
    все корни характеристического вещественны и различны.
\end{theorem}

\begin{proof}
    Если $\lambda_1, \dots, \lambda_n$ различны, то им соответствуют $n$ линейно независимых
    собственных векторов, следовательно, они составляют базис. А значит, матрица 
    $A$ диагональна.
\end{proof}

\begin{remark}
    Условие не является необходимым: в качестве примера можно привести тождественное
    преобразование.
\end{remark}

\begin{consequence}
    Пусть матрица преобразования $A$ диагональна и собственные значения вещественны.
    Тогда алгебраическая кратность любого собственного значения равна его геометрической 
    кратности.
\end{consequence}
\end{document}